% ==================================================================
% id:       app-012
% topic:    שימושי הנגזרת
% subtopic: מספר פתרונות של משוואה
% solution: full
% verified: no
% ==================================================================
\textbf{הוכיחו כי למשוואות הבאות קיים פתרון ממשי:}
\bE\bMs{2}
\item $\ds{ 2x^5 + 3x^3 + 6x - 10 = 0} $
\hint{השתמשו במשפט ערך הביניים: חפשו שתי נקודות שבהן לביטוי יש סימנים שונים.}
\begin{solbox}
\unverified
נסמן
$f(x)=2x^5 + 3x^3 + 6x - 10$.
זהו פולינום, ולכן
$f$
רציפה על כל הישר. מתקיים
$$f(0)=-10<0\qquad,\qquad f(1)=2+3+6-10=1>0$$
לכן לפי משפט ערך הביניים קיים
$c\in(0,1)$
עבורו
$f(c)=0$.

הערה: הפתרון אף יחיד, כי
$f'(x)=10x^4+9x^2+6>0$
לכל
$x$,
ולכן
$f$
עולה ממש ומתאפסת לכל היותר פעם אחת.
\end{solbox}
\item $\ds{1 + 4x + 3\cos x = 0}$
\begin{solbox}
\unverified
נסמן
$g(x)=1 + 4x + 3\cos x$,
פונקציה רציפה על כל הישר. מתקיים
$$g(0)=1+0+3=4>0\qquad,\qquad g(-\pi)=1-4\pi-3=-2-4\pi<0$$
לכן לפי משפט ערך הביניים קיים
$c\in(-\pi,0)$
עבורו
$g(c)=0$.

הערה: גם כאן הפתרון יחיד, כי
$g'(x)=4-3\sin x\geq 4-3=1>0$
לכל
$x$,
ולכן
$g$
עולה ממש.
\end{solbox}
\eMs
\eE
